\exo

\begin{enumerate}
\item Rappeler les critères de divisibilité par $2$, par $3$ et par $7$.
\item En utilisant l'égalité $6=2\times 3$, donner un critère de divisibilité par $6$.
\item Pour chacun des nombres suivants, indiquer s'il est divisible par $2$, par $3$, par $6$ et par $7$.
\[
42\qquad 1296\qquad 1029
\]
\end{enumerate}

\exo

Compléter le tableau suivant.

\medskip
\begin{center}
\begin{tabular}{|l|c|c|c|}
\hline
Nombre & Appartient à $\mathbb{N}$ & Appartient à $\mathbb{Z}$ & Écriture ensembliste \\
\hline
$2$ & Oui & Oui & $2\in\mathbb{N}\cap\mathbb{Z}$ \\
\hline
$-3$ & & & \\
\hline
$5$ & & & \\
\hline
$0$ & & & \\
\hline
$1000$ & & & \\
\hline
$-13$ & & & \\
\hline
\end{tabular}
\end{center}

\begin{enumerate}
\item Existe-t-il un nombre appartenant à $\mathbb{N}$ mais n'appartenant pas à $\mathbb{Z}$ ?
\item Existe-t-il un nombre appartenant à $\mathbb{Z}$ mais n'appartenant pas à $\mathbb{N}$ ?
\end{enumerate}

\exo

On considère un carré de côté $c$, avec $c\in\mathbb{N}$, et on note $A$ son aire.
\begin{enumerate}
\item Donner l'expression de $A$ en fonction de $c$.
\item Calculer $A$ pour $c=1,2,3,\ldots,10$.
\item Que peut-on dire de $A$ lorsque $c$ est pair ?
\item Que peut-on dire de $A$ lorsque $c$ est impair ?
\item Est-il possible d'avoir $A=7$ ? Même question pour $A=5$.
\item Si $A\leq 10$, lister toutes les valeurs possibles de $c$.
\end{enumerate}

\exo

On considère un rectangle de largeur $2$ et de longueur $L$, avec $L\in\mathbb{N}$.
On note $A$ son aire.
\begin{enumerate}
\item Exprimer $A$ en fonction de $L$.
\item L'aire $A$ peut-elle être impaire ?
\item Est-il possible d'avoir $A=27$ ? Si oui, déterminer $L$.
\item Est-il possible d'avoir $A=86$ ? Si oui, déterminer $L$.
\end{enumerate}

\exo

Dire si les affirmations suivantes sont vraies ou fausses. Justifier lorsque c'est faux.
\setlength{\columnseprule}{0pt}
\begin{multicols}{2}
\begin{enumerate}
\item Il n'existe aucun nombre pair multiple de $3$.
\item Si un nombre est multiple de $3$ et de $4$, alors il est multiple de $7$.
\item Pour tout $n\in\mathbb{N}$, le nombre $3n+9$ est un multiple de $3$.
\item $0$ est un nombre pair.
\item $0$ est un nombre premier.
\item $0$ est divisible par $3$.
\end{enumerate}
\end{multicols}

\exo

Recopier et compléter chaque phrase avec le mot \og multiple \fg{} ou le mot \og diviseur \fg{}.
\setlength{\columnseprule}{0pt}
\begin{multicols}{2}
\begin{enumerate}
\item $250$ est un \ldots{} de $50$.
\item $1000$ est un \ldots{} de $1$.
\item $30$ est un \ldots{} de $30$.
\item $5$ est un \ldots{} de $55$.
\item $3$ est un \ldots{} de $333\,666\,666$.
\end{enumerate}
\end{multicols}

\exo

\begin{enumerate}
\item Combien y a-t-il de multiples de $13$ entre $1$ et $1000$ ?
\item Combien y a-t-il de multiples de $13$ entre $1$ et $500$ ?
\item Combien y a-t-il de multiples de $13$ entre $500$ et $1000$ ?
\end{enumerate}

\exo

On note $\mathbb{P}$ l'ensemble des nombres premiers rangés dans l'ordre croissant.
On définit une fonction $\varphi$ sur $\mathbb{P}$ de la façon suivante : pour tout nombre premier $p$, $\varphi(p)$ est le rang de $p$ dans la liste des nombres premiers.

Par exemple, $\varphi(2)=1$ et $\varphi(3)=2$.

\begin{enumerate}
\item Déterminer $\varphi(11)$.
\item Déterminer le nombre premier $p$ tel que $\varphi(p)=10$.
\end{enumerate}
